\section{Montaje y metodología de pruebas}

Cada prueba se hizo en el entorno que mejor se podía controlar y se declara aquí junto con sus
límites. Las cifras de la sección de resultados salen de los registros que se guardaron como
evidencia, y cada una se puede rastrear hasta uno de esos archivos.

\subsection{Pruebas en el PC}

\begin{itemize}
  \item \textbf{Motor de cerca:} \dHostTests{} pruebas con \dHostChecks{} comprobaciones, y una prueba
        aleatoria de \dStressCases{} casos con semilla fija que compara la distancia del motor con
        la distancia de haversine. Se repiten bajo detectores de errores de memoria y de comportamiento
        indefinido.
  \item \textbf{Módulos de C puro} (\texttt{alert\_color}, \texttt{fence\_square}, \texttt{app\_fsm},
        \texttt{demo\_route}, \texttt{bench\_polygon}): cada uno tiene su prueba. Las de
        \texttt{fence\_square}, \texttt{demo\_route} y \texttt{bench\_polygon} usan el motor real y no
        una copia. La de \texttt{app\_fsm} recorre las 25 combinaciones de estado y evento.
  \item \textbf{Compilación cruzada:} el motor se compila para Cortex-M33 y ESP32-S3 con advertencias
        como errores. Compilar no demuestra que funcione en el ESP32-S3.
\end{itemize}

\subsection{Medida del tiempo de \texttt{geofence\_check()} en la placa}

Se compila el firmware con el benchmark activado. Al arrancar, antes de iniciar el GNSS y el botón,
construye polígonos regulares de 4, 16, 64, 256 y 1024 vértices con un radio de 500~m, y para cada uno
mide un lote de 40 llamadas: 20 puntos seguramente dentro y 20 seguramente fuera. Reporta el tiempo
total y la media por llamada, y cuenta las clasificaciones erróneas.

Como el contador del kernel corre a \dClockHz~Hz y cada lectura tiene \dTickUs~$\mu$s de resolución,
no se mide una llamada suelta sino el lote completo. Se repite en tres arranques consecutivos. El
tiempo es el del contador del kernel, no el de un osciloscopio sobre un pin.

\subsection{Modo demo}

Se pulsa el botón de la placa y se lee la consola mientras corre la ruta simulada. La ruta no usa
la posición real: sale de un centro ficticio. Se comprueba que el registro muestra el cambio de
estado, la creación de la cerca, los 42 pasos con su distancia y su color, y el regreso a la
búsqueda al pulsar de nuevo.

\subsection{Tiempo hasta el primer fix en frío}

El módem guarda datos de satélites en memoria no volátil, por lo que un reinicio normal es un
arranque en caliente. Para medir un arranque en frío se usa la compilación que borra esos datos en
cada arranque (se conserva solo el desfase del oscilador del propio dispositivo, como indica la
documentación de la librería del módem). La placa funciona a batería, en exteriores, sin PC. Con un
cronómetro se mide desde que se enciende hasta que el LED queda en verde fijo, que equivale al primer
fix con precisión reportada de 50~m o mejor. Se hicieron tres arranques; el tercero se dejó encendido
cinco minutos más para que el módem guardara sus datos.

Es una medida manual, con una incertidumbre de unos 2~s por la reacción de quien cronometra, y no
tiene registro de la consola.

\subsection{Precisión de la posición, prueba estática}

Con el firmware normal, que conserva los datos guardados, la placa se deja quieta, plana y boca
arriba, y se graba la consola. Se analiza el último arranque. Se descarta el primer minuto de fixes
porque el receptor aún converge, y se calcula la dispersión de las posiciones respecto a su propio
promedio, tras pasarlas a metros hacia el este y hacia el norte:

\begin{itemize}
  \item \textbf{CEP50:} la mediana de la distancia al promedio. La mitad de los fixes cae dentro de ese radio.
  \item \textbf{CEP95:} el percentil 95 de esa distancia.
  \item \textbf{DRMS:} la raíz del promedio de los cuadrados de la distancia.
  \item \textbf{Precisión reportada:} la estimación que da el propio receptor para cada fix (1 sigma, 2D).
\end{itemize}

Esto mide la \emph{precisión}, es decir cuánto se dispersan los fixes alrededor de su promedio. No
mide el error absoluto: no hubo una referencia de coordenadas conocidas ni un receptor RTK, y un error
constante que afecte a todos los fixes por igual no se ve. El cálculo lo hace un script que se
comprueba a sí mismo con puntos de resultado conocido y que solo escribe cifras derivadas, sin
coordenadas.

\subsection{Condiciones de las pruebas de campo}

Los tres arranques en frío se hicieron en una calle de barrio entre casas, con el cielo parcialmente
nublado, a media mañana. La prueba estática se hizo desde una ventana de primer piso orientada al sur, con la
ventana abierta y la placa prácticamente afuera, con un poco de árbol y cables eléctricos cerca. Ninguna
fue en campo abierto. No se registraron la altura de los obstáculos, el porcentaje de cielo visible ni
la interferencia de radio.
